\documentclass[10pt,a4paper]{amsart}
\usepackage{amsmath,amssymb,amsthm,graphicx}
\usepackage{verbatim}
\newenvironment{reference}{\comment}{\endcomment}
%\newenvironment{reference}{}{}
\newenvironment{slogan}{\comment}{\endcomment}
\newenvironment{history}{\comment}{\endcomment}

% For commutative diagrams we use Xy-pic
\usepackage[all]{xy}

% We use 2cell for 2-commutative diagrams.
\xyoption{2cell}
% Theorem environments.
%
\theoremstyle{plain}
\newtheorem{theorem}[subsection]{Theorem}
\newtheorem{proposition}[subsection]{Proposition}
\newtheorem{lemma}[subsection]{Lemma}

\theoremstyle{definition}
\newtheorem{definition}[subsection]{Definition}
\newtheorem{example}[subsection]{Example}
\newtheorem{exercise}[subsection]{Exercise}
\newtheorem{situation}[subsection]{Situation}

\theoremstyle{remark}
\newtheorem{remark}[subsection]{Remark}
\newtheorem{remarks}[subsection]{Remarks}

\numberwithin{equation}{subsection}

% Macros
%
\def\lim{\mathop{\mathrm{lim}}\nolimits}
\def\colim{\mathop{\mathrm{colim}}\nolimits}
\def\Spec{\mathop{\mathrm{Spec}}}
\def\Hom{\mathop{\mathrm{Hom}}\nolimits}
\def\Ext{\mathop{\mathrm{Ext}}\nolimits}
\def\SheafHom{\mathop{\mathcal{H}\!\mathit{om}}\nolimits}
\def\SheafExt{\mathop{\mathcal{E}\!\mathit{xt}}\nolimits}
\def\Sch{\mathit{Sch}}
\def\Mor{\mathop{\mathrm{Mor}}\nolimits}
\def\Ob{\mathop{\mathrm{Ob}}\nolimits}
\def\Sh{\mathop{\mathit{Sh}}\nolimits}
\def\NL{\mathop{N\!L}\nolimits}
\def\CH{\mathop{\mathrm{CH}}\nolimits}
\def\proetale{{pro\text{-}\acute{e}tale}}
\def\etale{{\acute{e}tale}}
\def\QCoh{\mathit{QCoh}}
\def\Ker{\mathop{\mathrm{Ker}}}
\def\Im{\mathop{\mathrm{Im}}}
\def\Coker{\mathop{\mathrm{Coker}}}
\def\Coim{\mathop{\mathrm{Coim}}}

% Boxtimes
%
\DeclareMathSymbol{\boxtimes}{\mathbin}{AMSa}{"02}

%
% Macros for moduli stacks/spaces
%
\def\QCohstack{\mathcal{QC}\!\mathit{oh}}
\def\Cohstack{\mathcal{C}\!\mathit{oh}}
\def\Spacesstack{\mathcal{S}\!\mathit{paces}}
\def\Quotfunctor{\mathrm{Quot}}
\def\Hilbfunctor{\mathrm{Hilb}}
\def\Curvesstack{\mathcal{C}\!\mathit{urves}}
\def\Polarizedstack{\mathcal{P}\!\mathit{olarized}}
\def\Complexesstack{\mathcal{C}\!\mathit{omplexes}}
% \Pic is the operator that assigns to X its picard group, usage \Pic(X)
% \Picardstack_{X/B} denotes the Picard stack of X over B
% \Picardfunctor_{X/B} denotes the Picard functor of X over B
\def\Pic{\mathop{\mathrm{Pic}}\nolimits}
\def\Picardstack{\mathcal{P}\!\mathit{ic}}
\def\Picardfunctor{\mathrm{Pic}}
\def\Deformationcategory{\mathcal{D}\!\mathit{ef}}
\usepackage{cleveref}
\crefname{originallemma}{original lemma}{original lemmas}
\Crefname{originallemma}{Original lemma}{Original lemmas}
\makeatletter
\newlabel{bootstrap-lemma-after-fppf-sep-lqf}{{046J}{source}{}{sourcepointer-1}{}}
\newlabel{bootstrap-lemma-after-fppf-sep-lqf@cref}{{[originallemma][0][]046J}{[1][0][]source}{}{}{}}
\makeatother
\crefname{originaldefinition}{original definition}{original definitions}
\Crefname{originaldefinition}{Original definition}{Original definitions}
\makeatletter
\newlabel{formal-spaces-definition-formal-algebraic-space}{{0AIM}{source}{}{sourcepointer-2}{}}
\newlabel{formal-spaces-definition-formal-algebraic-space@cref}{{[originaldefinition][0][]0AIM}{[1][0][]source}{}{}{}}
\makeatother
\crefname{originallemma}{original lemma}{original lemmas}
\Crefname{originallemma}{Original lemma}{Original lemmas}
\makeatletter
\newlabel{formal-spaces-lemma-presentation-representable}{{0AIG}{source}{}{sourcepointer-3}{}}
\newlabel{formal-spaces-lemma-presentation-representable@cref}{{[originallemma][0][]0AIG}{[1][0][]source}{}{}{}}
\makeatother
\crefname{originalreference}{original reference}{original references}
\Crefname{originalreference}{Original reference}{Original references}
\makeatletter
\newlabel{lemma-representable-flat}{{0GCG}{source}{}{sourcepointer-4}{}}
\newlabel{lemma-representable-flat@cref}{{[originalreference][0][]0GCG}{[1][0][]source}{}{}{}}
\makeatother
\crefname{originallemma}{original lemma}{original lemmas}
\Crefname{originallemma}{Original lemma}{Original lemmas}
\makeatletter
\newlabel{spaces-lemma-descent-representable-transformations-property}{{03KD}{source}{}{sourcepointer-5}{}}
\newlabel{spaces-lemma-descent-representable-transformations-property@cref}{{[originallemma][0][]03KD}{[1][0][]source}{}{}{}}
\makeatother
\crefname{originaltheorem}{original theorem}{original theorems}
\Crefname{originaltheorem}{Original theorem}{Original theorems}
\makeatletter
\newlabel{spaces-theorem-presentation}{{02WW}{source}{}{sourcepointer-6}{}}
\newlabel{spaces-theorem-presentation@cref}{{[originaltheorem][0][]02WW}{[1][0][]source}{}{}{}}
\makeatother
\title{The flat locus over a complete DVR is representable by a closed immersion of formal algebraic spaces}
\author{The Stacks Project Authors}
\date{stacks-project:restricted.tex:local-snapshot-20261012; GNU FDL 1.2 or later}
\begin{document}
\maketitle
\noindent\textit{Editorial note.} All delivered mathematical text is the author's own, copied from the pinned Stacks Project chapter file restricted.tex; no statement or proof step was rewritten, reordered or removed.\par\medskip
\noindent\textit{Editorial note.} The counted proof is the author's complete argument: an affine torsion-quotient construction, an etale-equivalence-relation gluing with monomorphism identifications, a descent of the closed immersion property, and a global universal property argument by etale covering, all inside the delivered span\par\medskip
\noindent\textit{Editorial note.} The label lemma-flat-locus names the counted statement, and the proof environment immediately following it in the source is delivered verbatim as its proof.\par\medskip
\noindent\textit{Editorial note.} The counted statement is carried by the author's own lemma environment in the chapter source.\par\medskip
\noindent\textit{Allowed standard prerequisites (editorial).}\begin{itemize}
\item Flatness over a complete DVR detected by vanishing of uniformizer-power torsion
\item The presentation definition of locally Noetherian formal algebraic spaces, with etale transition maps between affines
\item Descent of representability and of the closed immersion property along etale surjective morphisms, from the spaces and bootstrap lemmas cited
\end{itemize}\medskip
\ifdefined\hypertarget\hypertarget{sourcepointer-1}{}\fi
\noindent\textit{Source reference.} Original lemma 046J: sources/stacks-chow/tags.txt, source lines 5474--5474. Prerequisite pointer invoked at proof line(s) 5816; it names a standard fact of a sibling chapter that the counted argument uses, and every step of the counted proof built on it is retained verbatim in the delivered body.\par\medskip
\ifdefined\hypertarget\hypertarget{sourcepointer-2}{}\fi
\noindent\textit{Source reference.} Original definition 0AIM: sources/stacks-chow/tags.txt, source lines 13270--13270. Prerequisite pointer invoked at proof line(s) 5774; it names a standard fact of a sibling chapter that the counted argument uses, and every step of the counted proof built on it is retained verbatim in the delivered body.\par\medskip
\ifdefined\hypertarget\hypertarget{sourcepointer-3}{}\fi
\noindent\textit{Source reference.} Original lemma 0AIG: sources/stacks-chow/tags.txt, source lines 13264--13264. Prerequisite pointer invoked at proof line(s) 5790; it names a standard fact of a sibling chapter that the counted argument uses, and every step of the counted proof built on it is retained verbatim in the delivered body.\par\medskip
\ifdefined\hypertarget\hypertarget{sourcepointer-4}{}\fi
\noindent\textit{Source reference.} Original reference 0GCG: sources/stacks-restricted/pinned-original-source.tex, source lines 3136--3136. Prerequisite pointer invoked at proof line(s) 5793; it names a standard fact of the same chapter that the counted argument uses, and every step of the counted proof built on it is retained verbatim in the delivered body.\par\medskip
\ifdefined\hypertarget\hypertarget{sourcepointer-5}{}\fi
\noindent\textit{Source reference.} Original lemma 03KD: sources/stacks-chow/tags.txt, source lines 4673--4673. Prerequisite pointer invoked at proof line(s) 5818; it names a standard fact of a sibling chapter that the counted argument uses, and every step of the counted proof built on it is retained verbatim in the delivered body.\par\medskip
\ifdefined\hypertarget\hypertarget{sourcepointer-6}{}\fi
\noindent\textit{Source reference.} Original theorem 02WW: sources/stacks-chow/tags.txt, source lines 3819--3819. Prerequisite pointer invoked at proof line(s) 5805; it names a standard fact of a sibling chapter that the counted argument uses, and every step of the counted proof built on it is retained verbatim in the delivered body.\par\medskip
\begin{lemma}
\label{lemma-flat-locus}
Let $X$ be a locally Noetherian formal algebraic space over
a complete discrete valuation ring $A$.
Then there exists a closed immersion $X' \to X$
of formal algebraic spaces such that $X'$ is flat over $A$
and such that any morphism $Y \to X$ of locally Noetherian formal algebraic
spaces with $Y$ flat over $A$ factors through $X'$.
\end{lemma}

\begin{proof}
Let $\pi \in A$ be the uniformizer. Recall that an $A$-module
is flat if and only if the $\pi$-power torsion is $0$.

\medskip\noindent
First assume that $X$ is an affine formal algebraic space.
Then $X = \text{Spf}(B)$ with $B$ an adic Noetherian $A$-algebra.
In this case we set $X' = \text{Spf}(B')$ where
$B' = B/\pi\text{-power torsion}$. It is clear that $X'$ is flat
over $A$ and that $X' \to X$ is a closed immersion.
Let $g : Y \to X$ be a morphism of locally Noetherian formal algebraic spaces
with $Y$ flat over $A$. Choose a covering $\{Y_j \to Y\}$ as in
Formal Spaces, Definition
\ref{formal-spaces-definition-formal-algebraic-space}.
Then $Y_j = \text{Spf}(C_j)$ with $C_j$ flat over $A$.
Hence the morphism $Y_j \to X$, which correspond to a continuous
$R$-algebra map $B \to C_j$, factors through $X'$ as clearly
$B \to C_j$ kills the $\pi$-power torsion.
Since $\{Y_j \to Y\}$ is a covering and since $X' \to X$
is a monomorphism, we conclude that $g$ factors through $X'$.

\medskip\noindent
Let $X$ and $\{X_i \to X\}_{i \in I}$ be as in
Formal Spaces, Definition
\ref{formal-spaces-definition-formal-algebraic-space}.
For each $i$ let $X'_i \to X_i$ be the flat part as
constructed above. For $i, j \in I$ the projection
$X'_i \times_X X_j \to X'_i$ is an \'etale (by assumption)
morphism of schemes (by Formal Spaces, Lemma
\ref{formal-spaces-lemma-presentation-representable}).
Hence $X'_i \times_X X_j$ is flat over $A$ as morphisms
representable by algebraic spaces and \'etale
are flat (Lemma \ref{lemma-representable-flat}).
Thus the projection $X'_i \times_X X_j \to X_j$ factors
through $X'_j$ by the universal property. We conclude that
$$
R_{ij} = X'_i \times_X X_j = X'_i \times_X X'_j = X_i \times_X X'_j
$$
because the morphisms $X'_i \to X_i$ are injections of sheaves.
Set $U = \coprod X'_i$, set
$R = \coprod R_{ij}$, and denote $s, t : R \to U$ the two
projections. As a sheaf $R = U \times_X U$ and $s$ and $t$
are \'etale. Then $(t, s) : R \to U$ defines an \'etale equivalence
relation by our observations above. Thus $X' = U/R$ is an
algebraic space by Spaces, Theorem \ref{spaces-theorem-presentation}.
By construction the diagram
$$
\xymatrix{
\coprod X'_i \ar[r] \ar[d] & \coprod X_i \ar[d] \\
X' \ar[r] & X
}
$$
is cartesian. Since the right vertical arrow is \'etale surjective
and the top horizontal arrow is representable and a closed immersion
we conclude that $X' \to X$ is representable by
Bootstrap, Lemma \ref{bootstrap-lemma-after-fppf-sep-lqf}.
Then we can use Spaces, Lemma
\ref{spaces-lemma-descent-representable-transformations-property}
to conclude that $X' \to X$ is a closed immersion.

\medskip\noindent
Finally, suppose that $Y \to X$ is a morphism with
$Y$ a locally Noetherian formal algebraic space flat over $A$.
Then each $X_i \times_X Y$ is \'etale over $Y$ and
therefore flat over $A$ (see above).
Then $X_i \times_X Y \to X_i$ factors through $X'_i$.
Hence $Y \to X$ factors through $X'$ because
$\{X_i \times_X Y \to Y\}$ is an \'etale covering.
\end{proof}

\end{document}
